\begin{abstract}
Nuclear medicine is transitioning from a primarily qualitative diagnostic tool to a quantitative imaging science, but the partial volume effect (PVE) systematically undermines activity-concentration measurements in small or irregular lesions.
Standard recovery-coefficient (RC) corrections are calibrated on a handful of spherical inserts and do not transfer to the stochastic morphologies of clinical targets.
We present a proof of concept for a PVE validation framework based on radiolabeled tissue-equivalent gelatinous matrices.
The construction yields a high-density library of irregular, contrast-visible inserts: computed tomography (CT) supplies the morphological reference (volume and shape), the labeled gelatin preparation sets the activity density, and SPECT records the PVE-affected observation.
Across six experimental sessions we acquired $75$ SPECT/CT scans of $21$ phantom containers, yielding $2{,}302$ validated one-to-one CT--SPECT segment pairs in a cold (non-radioactive) gelatin background.
Activity density spanned $0.59$--$11.78\,\mu\mathrm{Ci/ml}$.
One session ($334$ segments) was held out as a new-preparation test.
The protocol is demonstrated with $^{99m}\mathrm{Tc}$ as a model tracer; the same construction is transferable to therapeutic isotopes such as $^{177}\mathrm{Lu}$ and to PET tracers.
Radiomic features extracted from the nuclear-medicine (NM) images were mapped to the reference targets with physics-anchored machine-learning models.
On the held-out session, a LightGBM model trained on the dimensionless ratio of activity density to mean NM counts recovered activity density with $R^2=0.95$, and a two-parameter affine calibration of NM volume recovered CT volume with $R^2=0.75$.
\end{abstract}

\section{Introduction}

Nuclear medicine is currently undergoing a transformative paradigm shift, transitioning from a primarily qualitative diagnostic tool to a robust quantitative imaging science.
In the contemporary era of personalized oncology and Theranostics, the ability to accurately measure absolute radiotracer uptake is no longer an auxiliary feature but a clinical imperative~\citep{soret2007partial,erlandsson2012review}.
High-stakes therapeutic decisions, such as determining patient eligibility for targeted radionuclide therapies like $^{177}\mathrm{Lu}$-PSMA, are increasingly governed by specific Standardized Uptake Value (SUV) thresholds derived from PET imaging~\citep{jentzen2011phantom}.
Furthermore, the rising importance of internal dosimetry and the utilization of bone scintigraphy as a quantitative biomarker demand a level of precision that transcends traditional visual interpretation~\citep{zeintl2010quantitative,seret2012quantitative}.

However, the clinical reliability of these quantitative metrics is fundamentally undermined by the Partial Volume Effect (PVE).
PVE is a pervasive phenomenon in emission tomography---affecting both Positron Emission Tomography (PET) and Single-Photon Emission Computed Tomography (SPECT)---that leads to a significant degradation of signal fidelity~\citep{soret2007partial,erlandsson2012review,aston2002positron}.
The clinical consequence of PVE is most severe in small lesions or regions characterized by high activity gradients, where it results in a systematic underestimation of activity concentration~\citep{hoffman1979quantitation,kessler1984analysis}.
When a small malignant focus is subject to these effects, its measured SUV may fall below critical clinical cut-off points, potentially leading to false-negative diagnoses or the denial of life-saving treatments to eligible patients.

Technically, PVE is driven by the finite spatial resolution of imaging systems and is commonly described through two primary mechanisms.
The first is the system's Point-Spread Function (PSF), which causes a ``spill-over'' (cross-contamination) of counts between adjacent regions~\citep{aston2002positron,rousset1998correction,frouin2002correction}.
This manifests as spill-out from the target lesion into the background and spill-in from surrounding tissues into the region of interest.
The second mechanism is the tissue fraction effect, where a single voxel encapsulates a mixture of tissues with varying uptake levels, further blurring the boundaries and complicating accurate quantification~\citep{meltzer1990correction,mullergartner1992measurement}.

To mitigate these inaccuracies, various correction strategies have been proposed over the past decades, including recovery coefficient (RC) methods, deconvolution-based algorithms, and structural image-guided approaches~\citep{rousset1998correction,baete2004evaluation}.
While showing promise, many of these techniques suffer from high sensitivity to noise or a heavy reliance on perfect anatomical segmentation~\citep{erlandsson2012review,hoetjes2010partial}.
Historically, validation studies have utilized phantoms containing a limited number of spherical inserts of known sizes to calculate RCs~\citep{hoffman1979quantitation,geworski2000recovery}.
However, more recent research has demonstrated that geometry plays a decisive role; for instance, circular and elliptical structures of identical volumes exhibit disparate recovery coefficients~\citep{zeintl2010quantitative,gallivanone2011pve}.
This implies that clinical lesions---which are inherently irregular and stochastic in shape---possess unpredictable RCs that cannot be accurately corrected using idealized, rigid geometric models~\citep{jentzen2011phantom,seret2012quantitative}.

In this study, we present a Proof of Concept for a novel PVE validation framework based on radiolabeled tissue-equivalent gelatinous matrices.
Unlike traditional ``perfect'' geometric phantoms, the unique physical properties of this medium allow for the efficient generation of a vast repository of stochastic, asymmetrical, and morphologically heterogeneous segments.
By leveraging this method, we successfully produced a high-density dataset of $2{,}302$ unique validated segments, each with a CT-derived volume and a preparation-defined activity density.

This dataset allows for an advanced analysis of PVE that was previously unattainable with sparse phantom data.
Because the irregular geometry of these segments causes traditional mathematical recovery models to collapse, we employed Radiomics to capture a high-dimensional feature set, including complex spatial and textural patterns that are often ``smoothed'' by PVE.
These features were then integrated into a Machine Learning (ML) framework to bridge the statistical gap between image-derived values and ground-truth metrics.
This synergy of morphological realism and predictive modeling aims to provide a robust, data-driven PVE correction that is uniquely adaptable to the intricate irregularities of real-world clinical findings~\citep{salmon2018quantitative,trangia2019characterization,southekal2012evaluation,matsubara2020model,denijs2014improving}.

\section{Methods}

\subsection{General algorithm and experimental rationale}

The core objective of this methodology is to provide a data-driven solution to the Partial Volume Effect (PVE) by creating a high-throughput data factory for SPECT/CT analysis.
The algorithm is based on a four-stage pipeline designed to bridge the gap between idealized phantoms and clinical reality:

\paragraph{Contrast-based ground truth.}
We developed a two-phase phantom system where the background matrix consists of non-active gelatin, and the target segments consist of gelatin radiolabeled with $^{99m}\mathrm{Tc}$ and enriched with an iodinated CT contrast agent.
This design allows the CT scan to serve as the morphological ground truth, providing volume and shape data, while the SPECT scan captures the distorted functional data affected by PVE.

\paragraph{Stochastic data augmentation.}
To ensure statistical robustness, multiple containers were utilized.
By physically agitating (shaking) the containers between scans and utilizing the natural decay of the isotope across multiple experimental sessions, we generated $2{,}302$ unique, non-idealized segments from a single experimental framework.

\paragraph{Proof of concept and scalability.}
$^{99m}\mathrm{Tc}$ was strategically selected as a model isotope due to its availability and convenient half-life, enabling the rapid generation of a large dataset.
This establishes a cost-effective methodology that is fundamentally generalizable to more expensive or logistically complex isotopes, such as $^{177}\mathrm{Lu}$ or various PET tracers.

\subsection{Phantom preparation and formulation}

To maintain physical consistency and ensure that the only variables were activity and contrast, a standardized preparation protocol was followed.

\paragraph{Base gelatin matrix.}
Distilled water was heated to a constant $65^\circ\mathrm{C}$ using a precision heating bath.
Gelatin powder was integrated at a ratio of $19\,\mathrm{g/L}$ and stirred for approximately one hour until a uniform, fully melted distribution was achieved.
This identical base matrix was used for both the background and the active segments to ensure matching attenuation and density properties of the gelatin itself.

\paragraph{Active segments formulation.}
For the target segments, a dark-pigmented juice was used as a solvent to aid visual identification during preparation (\Cref{fig:phantom}, right).
An iodinated contrast agent ($10\,\mathrm{mL/L}$) was added to ensure high-contrast visibility and facilitate automated, high-precision segmentation in the CT scan (\Cref{fig:phantom}, left).

\paragraph{Radiolabeling.}
Radioactive $^{99m}\mathrm{Tc}$ was integrated into these segments at varying activity densities across different sessions.
These segments were then randomly distributed within the non-active background gelatin inside multiple containers.

\subsection{Data volume and acquisition density}

The high-throughput nature of the experiment is summarized in \Cref{tab:sessions}.
Data were collected over six experimental sessions, generating a total of $2{,}302$ validated segments.

\paragraph{Dynamic range.}
The activity density across the sessions spanned from $0.59$ to $11.78\,\mu\mathrm{Ci/ml}$, providing the models with a wide range of signal intensities.

\paragraph{Acquisition density.}
A total of $75$ scans were performed across $21$ unique phantom containers.
Between scans, the phantoms were physically agitated to redistribute the segments, and scans were repeated over time to utilize natural radioactive decay, thereby increasing the diversity of the feature space.

\paragraph{Test dataset integrity.}
Consistent with our validation strategy, session~$0$ containing $334$ segments was held out entirely for independent testing (``new phantom'' simulation).
This test set was characterized by two disjoint activity ranges ($3.66$--$5.23$ and $7.15$--$10.23\,\mu\mathrm{Ci/ml}$) that lie inside the global training envelope but form a different mix, so that accuracy is not guaranteed by memorizing the training distribution.

\begin{table}[ht]
\centering
\caption{Summary of experimental data acquisition.
Sessions $1$--$5$ served as the training/validation pool; session~$0$ is the independent hold-out.
Segment counts, scan counts and activity ranges are recomputed from the released feature tables (\texttt{scripts/make\_paper\_figures.py}); jar counts are from the experimental log.}
\label{tab:sessions}
\small
\begin{tabular}{lrrrrl}
\toprule
Session & Segments & Scans & Jars & Activity ($\mu$Ci/ml) \\
\midrule
$1$ & $386$ & $11$ & $3$ & $5.5$--$7.3$ \\
$2$ & $327$ & $13$ & $3$ & $3.4$--$4.7$ \\
$3$ & $190$ & $10$ & $3$ & $3.04$--$3.96$ \\
$4$ & $583$ & $15$ & $4$ & $0.59$--$0.90$ \\
$5$ & $482$ & $13$ & $4$ & $1.7$--$2.5$; $8.4$--$11.8$ \\
$0$ (held-out) & $334$ & $13$ & $4$ & $3.6$--$5.2$; $7.1$--$10.2$ \\
\bottomrule
\end{tabular}
\end{table}

\subsection{Imaging protocol and reconstruction}

Data acquisition was performed using a Siemens Symbia Pro.specta SPECT/CT scanner at Shaare Zedek Medical Center.
The SPECT protocol consisted of $360^\circ$ projections with an angular dwell time of $3$ seconds per degree.
Images were reconstructed using the Ordered Subset Expectation Maximization (OSEM) algorithm with attenuation correction (AC), $48$ iterations, and $1$ subset.
A clinical-standard $3$D Gaussian filter ($8\,\mathrm{mm}$ FWHM) was applied to the final volumes to simulate standard clinical resolution.

\subsection{Automated segmentation and data integrity}

A dedicated multi-stage pipeline was developed to match functional and anatomical data (\Cref{fig:segmentation}).

\paragraph{CT segmentation (ground truth).}
Anatomical boundaries were defined using a fixed threshold of $60\,\mathrm{HU}$.

\paragraph{NM segmentation.}
Functional volumes were initially thresholded at $8\%$ of the global maximum, followed by a local refinement at $40\%$ of the segment's peak voxel value.

\paragraph{Recursive refinement.}
If a refined segment's volume was $<60\%$ of its initial estimate, the algorithm recursively searched for a secondary local maximum ($>70\%$ of the primary peak) and reapplied the refinement.
This iterative approach is intended to keep segments representative of irregular, non-spherical morphologies when two inserts have merged on SPECT.

\paragraph{Data cleaning.}
To ensure a unique correspondence between the morphological reference and the observation, any segment without a strict $1$:$1$ CT-to-NM match (for example overlapping segments or merged clusters) was excluded from the final dataset.

\begin{figure}[t]
\centering
\begin{minipage}[t]{0.48\linewidth}
\centering
\includegraphics[width=\linewidth]{fig1_ct.png}
\end{minipage}\hfill
\begin{minipage}[t]{0.48\linewidth}
\centering
\includegraphics[width=\linewidth]{fig1_photo.png}
\end{minipage}
\caption{Phantom structure overview.
Left: volume-rendered CT scan, showing the 3D spatial distribution of contrast-enhanced segments (red) and the background matrix (blue).
This data is used to define the morphological ground truth.
Right: corresponding photograph of the physical stochastic phantom containers, illustrating the appearance of the embedded irregular gelatin segments.}
\label{fig:phantom}
\end{figure}

\begin{figure}[t]
\centering
\includegraphics[width=0.92\linewidth]{fig2_segmentation.png}
\caption{Phantom segmentation comparison.
Left: top row shows the PVE-affected NM scan; bottom row shows the NM mask.
Right: top row shows the corresponding high-resolution CT scan; bottom row shows the precise CT mask.
This comparison highlights the blurring and resolution loss in functional data, demonstrating how CT-based masks define the morphological ground truth for validating PVE correction models.}
\label{fig:segmentation}
\end{figure}

\subsection{Radiomics and machine learning architecture}

The experimental library comprises $2{,}302$ validated one-to-one segments.
Predictive models were trained on the reduced NM-40 radiomics table released with this repository ($1{,}766$ training segments / $315$ held-out segments, $79$ features): the smallest ROIs are already dropped, and every CT-derived column is excluded because several of those columns are algebraic restatements of the CT volume target.

\paragraph{Feature extraction and selection.}
High-dimensional radiomic features were extracted from the NM data using a fixed bin width of $2{,}000$ counts (the value recorded in the PyRadiomics settings column of every row).
This standardization keeps a constant intensity resolution across activity levels and scan sessions.
To prevent overfitting and reduce collinearity, constant features and pairs with absolute Pearson correlation above $0.995$ were excluded.
Absolute position inside the phantom (centre of mass, bounding-box origin) is also removed, so that a model cannot memorise insert layout.

\paragraph{Predictive modeling strategy.}
The two targets are treated separately.

Activity density is proportional to the mean reconstructed intensity $\bar{c}$ in the NM mask.
The baseline is therefore the one-parameter calibration
\begin{equation}
\label{eq:activity-cal}
A = k\,\bar{c},
\end{equation}
with $k$ fitted through the origin on the training sessions.
Radiomic features enter only as a correction to this scale.
A gradient-boosting model (LightGBM) is trained to predict the ratio $A/\bar{c}$ from the NM feature vector; the activity-density estimate is that predicted ratio multiplied by the observed $\bar{c}$.
A LightGBM model trained to predict $A$ directly, rather than the ratio, is reported as a comparison.

NM-mask volume $V_{\mathrm{NM}}$ is linearly related to CT volume, but the two are not proportional.
The volume estimate is the two-parameter fit
\begin{equation}
\label{eq:volume-cal}
V_{\mathrm{CT}} = a\,V_{\mathrm{NM}} + b,
\end{equation}
with $a$ and $b$ estimated on the training sessions.
An elastic-net regressor on the same NM features is included only as a comparison.

\paragraph{Validation.}
Sessions~$1$--$5$ are used for training.
Cross-validation holds out one acquisition date at a time.
Reported scores are computed on session~$0$, which is not used for training or selection.
Activity density is prepared once per acquisition and is constant, or nearly constant, within a scan ($51$ of $67$ training scans have a single value).
Segments from the same scan are therefore not independent, and a random split of segments is not used to select or compare models.

\section{Results}

All scores below are computed on session~$0$ of the reduced NM-40 table: $315$ segments from $13$ acquisitions that share no scan and no date with training ($1{,}766$ segments, $67$ acquisitions, $79$ features).
On this session both targets lie above the training mean (activity density $6.76$ versus $4.08\,\mu\mathrm{Ci/ml}$; CT volume $19.2$ versus $15.4\,\mathrm{ml}$).
\Cref{tab:ml} reports session-grouped cross-validation on the training dates together with the held-out scores.

\paragraph{Activity density.}
Mean reconstructed intensity $\bar{c}$ is linearly related to prepared activity density (\Cref{fig:activity}, left).
The held-out session occupies two disjoint bands ($3.6$--$5.2$ and $7.1$--$10.2\,\mu\mathrm{Ci/ml}$) that fall on the same line as the training sessions.
The horizontal banding in that panel is the per-acquisition preparation: many scans contribute a single activity value.
The global line is nevertheless dominated by between-session differences in prepared activity; within a session the correlation of $\bar{c}$ with activity density falls to approximately $0.45$--$0.55$ for half of the training sessions.

The through-origin calibration~\eqref{eq:activity-cal} already recovers the held-out densities with $R^2=0.93$, $\mathrm{RMSE}=0.61\,\mu\mathrm{Ci/ml}$ and $\mathrm{MAPE}=7.2\%$ (session-grouped CV $R^2=0.92$).
The observed ratio $A/\bar{c}$ is stable across sessions, including session~$0$ (approximately $2.8$--$3.4\times 10^{-5}\,\mu\mathrm{Ci\,ml^{-1}}$ per count).

LightGBM trained on $A/\bar{c}$ improves these scores to $R^2=0.95$, $\mathrm{RMSE}=0.51\,\mu\mathrm{Ci/ml}$, $\mathrm{MAE}=0.41\,\mu\mathrm{Ci/ml}$ and $\mathrm{MAPE}=6.7\%$ (session-grouped CV $R^2=0.93$).
The identity plot (\Cref{fig:activity}, right) follows the two held-out bands along the line of equality, with a small overestimate in the lower band and larger scatter in the upper band.
A LightGBM model trained to predict $A$ directly reaches only $R^2=0.43$ ($\mathrm{RMSE}=1.72\,\mu\mathrm{Ci/ml}$, $\mathrm{MAPE}=23.8\%$).
The radiomic step is therefore a modest correction to a calibration that already transfers to a new preparation.

\begin{figure}[ht]
\centering
\includegraphics[width=0.98\linewidth]{fig3_activity.png}
\caption{Activity density on the reduced NM-40 table.
Left: prepared activity density versus mean NM intensity for the training sessions (blue) and session~$0$ (red).
The held-out points occupy two activity bands on the same line as training.
Right: predicted versus true activity density on session~$0$ for LightGBM trained on $A/\bar{c}$.
The dashed line is identity.}
\label{fig:activity}
\end{figure}

\paragraph{Volume.}
CT volume is an affine function of the $40\%$ NM-mask volume (\Cref{fig:volume}, left).
The fit on the training sessions is $V_{\mathrm{CT}} \approx 0.893\,V_{\mathrm{NM}} + 6.00$.
Held-out volumes ($6.8$--$51.4\,\mathrm{ml}$) overlap the main training cloud ($1.9$--$71.8\,\mathrm{ml}$).

On session~$0$ this fit attains $R^2=0.75$, $\mathrm{RMSE}=3.78\,\mathrm{ml}$, $\mathrm{MAE}=2.80\,\mathrm{ml}$ and $\mathrm{MAPE}=14.5\%$ (session-grouped CV $R^2=0.70$).
The identity plot (\Cref{fig:volume}, right) is tight between about $10$ and $25\,\mathrm{ml}$; the largest inserts are under-predicted.
An elastic-net regressor on the same $79$ NM features does not improve on the affine fit (held-out $R^2=0.67$, $\mathrm{RMSE}=4.30\,\mathrm{ml}$, $\mathrm{MAPE}=18.5\%$; session-grouped CV $R^2 < 0$).
Under SPECT-only features, the NM-to-CT volume map is limited by that physical relationship rather than by the choice of regressor.

\begin{figure}[ht]
\centering
\includegraphics[width=0.98\linewidth]{fig4_volume.png}
\caption{Volume on the reduced NM-40 table.
Left: CT volume versus $40\%$ NM-mask volume for the training sessions (blue) and session~$0$ (red).
Right: predicted versus true CT volume on session~$0$ for the affine calibration~\eqref{eq:volume-cal}.
The dashed line is identity.}
\label{fig:volume}
\end{figure}

\begin{table}[ht]
\centering
\caption{Held-out performance on session~$0$ of the reduced NM-40 table, with session-grouped cross-validation on the training dates.
Activity errors are in $\mu$Ci/ml; volume errors are in ml.}
\label{tab:ml}
\small
\begin{tabular}{llcccc}
\toprule
Target & Model & Session-CV $R^2$ & Test $R^2$ & Test RMSE & Test MAPE \\
\midrule
Activity & LightGBM on $A$ & $0.58$ & $0.43$ & $1.72$ & $23.8\%$ \\
Activity & $A = k\,\bar{c}$ & $0.92$ & $0.93$ & $0.61$ & $7.2\%$ \\
Activity & LightGBM on $A/\bar{c}$ & $0.93$ & $0.95$ & $0.51$ & $6.7\%$ \\
Volume & Elastic net & ${<}\,0$ & $0.67$ & $4.30$ & $18.5\%$ \\
Volume & $V_{\mathrm{CT}} = a V_{\mathrm{NM}}+b$ & $0.70$ & $0.75$ & $3.78$ & $14.5\%$ \\
\bottomrule
\end{tabular}
\end{table}

\FloatBarrier

\section{Summary}

This paper presented a method for estimating the activity density and the volume of irregular gelatin inserts from SPECT radiomics and the corresponding NM measurements.
The experimental library is a two-compartment gelatin phantom imaged with SPECT/CT, producing $75$ scans of $21$ containers over six sessions and $2{,}302$ validated one-to-one pairs, with session~$0$ held out as a new-preparation test.
For activity density, the count calibration~\eqref{eq:activity-cal} already reaches $R^2=0.93$ on the held-out session; LightGBM trained on the ratio $A/\bar{c}$ improves this to $R^2=0.95$.
For volume, the affine calibration~\eqref{eq:volume-cal} reaches $R^2=0.75$; radiomic regression does not improve on that fit.

The accompanying code (\texttt{src/phantom\_jelly/} and \texttt{scripts/}) regenerates the session inventory and \Cref{fig:activity,fig:volume} from the released CSVs.

\bibliographystyle{plainnat}
\bibliography{references}
